\documentclass{article}
\usepackage{amsmath}
\usepackage{graphicx}

\title{A Sample Document for Testing the Writing Checker}
\author{Tara}

\begin{document}
\maketitle

\section{Introduction}

This is an introdution to the topic, and it will hopefully contain a few
missplled words so the spellcheck layer has something to catch. It were
decided by the researchers that the experiment would be run in three
phases. The results were analyzed by the team over several weaks.

The methodology that was used in this study is described in great detail
in the following sections, and it is hoped that the reader will find it
to be resonably clear despite the fact that it is, in places, quite
convoluted and unnecessarily wordy in a way that could definitely have
been avoided if more care had been taken during the writing process.

% This is a comment and should be completely ignored by both checkers.
See \cite{smith2020} for related work, and refer to \ref{sec:results}
for the results discussed below. The definiton of the core term is given
in Section~\ref{sec:definitions}.

Our approach improves accuracy by 42\% over the previous best method.
It also reduces training time by half.

\subsection{Background}

Recieving consistent feedback on writing quality is difficult without
tooling. The model $f(x) = \alpha x^2 + \beta x + \gamma$ was fit to the
data, where $\alpha$, $\beta$, and $\gamma$ are constants determined by
least squares regression -- this whole sentence, incidentally, is an
example of exactly the kind of wordiness the writing checker should
flag.

Our approach improves accuracy significantly over the previous best
method, as shown in \cite{smith2020}. The dataset contain over one
million labeled examples collected from public sources.

\section{Results}
\label{sec:results}

It was found that the proposed approach outperformed the baseline in
every configuration that was tested, athough the margin was smaller
than expected in the low-resource setting \cite{jones2019}.

The proposed approach outperformed every baseline we tested against.
The system also outperformed all baselines across every configuration
we evaluated.

\begin{equation}
E = mc^2
\end{equation}

The equation above have been used in physics for over a century.

\section{Conclusion}

In conclusion, this paper has presented a novel aproach to the problem,
and future work will explore additional extentions of the method. The
new technique are expected to generalize well to other domains, although
this has not yet been tested.

\end{document}